% project_report.tex — ERK-format report template for YOUR project (English).
% Fill every [PLACEHOLDER]. Same two-column style as the real ERK proceedings.
% Compile like any deck: pdflatex project_report.tex (erk.sty sits next to it).
\documentclass[a4paper]{article}

\usepackage[utf8]{inputenc}\usepackage[T1]{fontenc}
\usepackage{erk}
\usepackage{times}
\usepackage{graphicx}
\usepackage[english]{babel}
\usepackage[top=22.5mm, bottom=22.5mm, left=22.5mm, right=22.5mm]{geometry}
\usepackage{booktabs,hyperref,array}

% keep the ERK look: no footnotes on the cover page
\def\footnotemark{}

\begin{document}

\title{flagwire: what a local 27B model can build against a locked specification}

\author{Bojan Jak\v{s}i\'{c}$^{1}$}

\affiliation{$^{1}$University of Ljubljana, Faculty of Computer and Information Science}

\email{Email: [your@student.uni-lj.si]}

\maketitle

\begin{abstract}{flagwire: what a local 27B model can build against a locked specification}
[English abstract, 200 words at most: what the software does, who it is for,
what the agents did, and the one result you can show. Write it last, readers
read it first.]
\end{abstract}

\section{Introduction}
What problem does your project solve, and who is the user? Why is it worth
building? Two short paragraphs maximum.

\section{Goals and objectives}
Numbered, checkable. For each goal someone can tell whether you met it.
\begin{itemize}
% Copied from README.md, "Goals"; report each as met or not, with the evidence.
\item[G1.] All 31 tier 1 acceptance tests pass.
\item[G2.] At least 20 of the 25 tier 2 acceptance tests pass.
\item[G3.] On Glas v0.1.0, \texttt{flagwire apply admin} gates its six admin routes, Glas's tests pass, and a pull request is opened.
\item[G4.] In the demo, \texttt{flagwire toggle} switches Glas's admin pages off and on without a restart.
\item[G5.] At least 80\% of the commits that change \texttt{src/} are the local model's; every other change is a logged intervention.
\item[G6.] Pass rate per tier and level over time, interventions per slice, and tokens and time per session are reported and reproducible.
\item[G7.] Both UX test runs are recorded; tier 1 averages a clarity of at least 3.5 of 5.
\end{itemize}

\section{Agentic stack used}
What you actually ran, pinned, with what each part did. Honest table, not
marketing. Model tags are part of the method --- pin them.

\begin{table}[h]
\centering\small
\begin{center}
\begin{tabular}{ | >{\raggedright\arraybackslash}p{0.22\linewidth} | >{\raggedright\arraybackslash}p{0.36\linewidth} | >{\raggedright\arraybackslash}p{0.30\linewidth} | }
\hline 
  \textbf{Tool} & \textbf{Model (pinned tag)} & \textbf{What it did}\\
\hline
  % From docs/STACK.md; fill the TBD pins.
  OpenCode [TBD] & Qwen3.8 27B GSQ RCO IQ3\_S (LM Studio [TBD]) & wrote \texttt{src/}, one slice per session \\
  Claude Code (web) & Claude Opus 5.5 (\texttt{claude-opus-5-5}) & spec, locked tests, plan, reviews \\
  flagwire \texttt{propose} & the same local model & groups routes and toggles into features \\
\hline
\end{tabular}
\end{center}
\caption{The stack as used, not as advertised.}
\end{table}

\section{Method --- how you directed the agents}
How work was split, what files carried state (specs, progress files), how you
verified each agent output before committing. Your commit log is the
replication record: which prompt and which model produced which part.

\section{Results}
What works today; how you know (tests you ran, the demo, the counts). What
failed and stayed broken is a result too --- report it.

\section{Limitations and next steps}
What this version does not do; what you would do next; what you would change
about the process, not just the product.

\section{AI-use declaration}
What the agents produced versus what you produced; every model with its exact
pinned tag and date of use; what was auto-generated at release (model defaults).
This section is graded as part of the report; write it like a method section.

\small
\begin{thebibliography}{1}
\bibitem{repo} Repository: \url{https://github.com/bjaksic84/bjaksic84-acs-project}, tag [TBD] used for the demo.
\end{thebibliography}

\end{document}
